\documentclass[11pt]{article}
\usepackage[margin=0.88in]{geometry}
\usepackage{amsmath,amssymb,amsthm}
\usepackage{graphicx}
\usepackage{hyperref}
\usepackage{microtype}
\hypersetup{hidelinks}

\newtheorem{theorem}{Theorem}
\newtheorem{lemma}[theorem]{Lemma}
\newcommand{\C}{\mathbb C}
\newcommand{\R}{\mathbb R}
\newcommand{\Rea}{\operatorname{Re}}
\newcommand{\Ima}{\operatorname{Im}}

\title{A Planar Counterexample to Absolute\
Leray-Kernel Boundedness on Strongly $\C$-Convex Domains}
\author{Candidate counterexample packet for arXiv:2506.02913}
\date{11 August 2026}

\begin{document}
\maketitle

\begin{center}
\fbox{\parbox{0.92\textwidth}{\textbf{Status.}
Candidate full negative answer to Remark~7.3 as literally stated; likely
valid, pending human review.  The question quantifies over $D\subset\C^n$
without imposing $n\ge2$.  In dimension one, strong $\C$-convexity is
vacuous beyond ordinary $\C$-convexity, and an explicit smooth radial domain
has an interior zero of the denominator $B_D$.  The positive kernel is then
nonintegrable on an open set.  If the intended question silently assumes
$n\ge2$, this packet does not settle that strengthened version.}}
\end{center}

\section{Source question}

Chatterjee, \emph{Dual realizations of Bergman spaces on strongly convex
domains}, arXiv:2506.02913v2 [1], proves $L^p$ boundedness of the positive
operator $|\mathcal B_D|$ when $D$ is strongly real convex.  Remark~7.3 asks:

\begin{quote}
Let $D\subset\C^n$ be a bounded $C^2$-smooth strongly $\C$-convex domain
containing the origin, and let $|\mathcal B_D|$ be the operator in
Lemma~3.1.  Is $|\mathcal B_D|$ an $L^p(D)$-bounded operator?
\end{quote}

\begin{figure}[h]
\centering
\includegraphics[width=0.94\textwidth]{figures/open_problem_crop.png}
\caption{Remark~7.3 and the exact question on source PDF page~27.}
\end{figure}

For a star-shaped domain with Minkowski functional $m_D$, the source defines
\begin{align}
 B_D(\zeta,z)
 &=m_D(\zeta)\,2\partial m_D(\zeta)(\zeta-z)
   +1-m_D(\zeta)^2,                                      \label{eq:B}\\
 |\mathcal B_D|f(z)
 &=\int_D f(\zeta)|B_D(\zeta,z)|^{-n-1}\,dV(\zeta).
                                                               \label{eq:op}
\end{align}
The pairing in \eqref{eq:B} is ordinary multiplication when $n=1$.

\section{Proof intuition}

The strongly real-convex proof uses the sign of the real supporting
hyperplane to keep $B_D$ away from zero in $D\times D$.  Complex convexity
controls only the modulus of a boundary complex support function; it does not
give that sign.  In one complex dimension there are no nonzero complex
tangent directions at all.  We therefore choose a smooth two-lobed radial
domain.  At a carefully selected point on one inward-facing side, the normal
ray re-enters the opposite lobe.  Moving partway down the original radial
segment and then along that normal gives a second interior point at which
$B_D$ vanishes exactly.

The zero is transverse.  Hence the implicit-function theorem supplies an
interior zero in the $\zeta$ variable for every $z$ in an open neighborhood.
Since $n+1=2$, the local kernel dominates $|\zeta-\zeta(z)|^{-2}$, whose area
integral diverges logarithmically.

\section{The domain and its gauge}

For $\theta\in\R$ set
\begin{equation}\label{eq:R}
 R(\theta)=1+\frac{13}{20}\cos(2\theta),
 \qquad
 D=\{re^{i\theta}:0\le r<R(\theta)\}.
\end{equation}
Because $R\ge7/20$, this is a bounded $C^\infty$ radial domain containing
$0$, star-shaped with respect to $0$, with only non-radial tangents.  It is
simply connected, hence $\C$-convex in dimension one.  Moreover, for any
$w\in\partial D$, the complex tangent space
\[
 H_w\partial D=\{v\in\C:\partial\rho(w)v=0\}
\]
is $\{0\}$.  Thus the usual condition that the second fundamental form be
positive definite on $H_w\partial D$ holds vacuously.  Accordingly, $D$ is
strongly $\C$-convex under the standard dimension-one convention used by the
question as written.

Its Minkowski functional is
\begin{equation}\label{eq:gauge}
 m_D(re^{i\theta})=\frac r{R(\theta)}.
\end{equation}
If $\eta=R(\theta)e^{i\theta}\in\partial D$, polar differentiation gives
\begin{equation}\label{eq:q}
 q(\theta):=2\partial m_D(\eta)
 =e^{-i\theta}\left(\frac1{R(\theta)}
             +i\frac{R'(\theta)}{R(\theta)^2}\right).
\end{equation}
Positive homogeneity also gives
$2\partial m_D(s\eta)=q(\theta)$ for every $s>0$.

\section{An exact transverse interior zero}

\begin{theorem}\label{thm:main}
For the domain \eqref{eq:R}, $|\mathcal B_D|$ fails to map $L^p(D)$ into
finite-valued functions for every $1\le p\le\infty$.  In particular it is
not bounded on $L^p(D)$ for any $1<p<\infty$.
\end{theorem}

\begin{proof}
Choose $\theta_0$ in the second quadrant by
\begin{equation}\label{eq:theta}
 \cos\theta_0=-\frac15,
 \qquad \sin\theta_0=\frac{2\sqrt6}{5}.
\end{equation}
Then
\[
 \cos(2\theta_0)=-\frac{23}{25},\qquad
 \sin(2\theta_0)=-\frac{4\sqrt6}{25},\qquad
 R(\theta_0)=\frac{201}{500}.
\]
Put $s=1/2$, $\eta_0=R(\theta_0)e^{i\theta_0}$,
$\zeta_0=s\eta_0$, and abbreviate $q_0=q(\theta_0)$.  Define
\begin{equation}\label{eq:zdef}
 z_0=\zeta_0+\frac{1-s^2}{s q_0}.
\end{equation}
Both points are nonzero, and $m_D(\zeta_0)=s<1$.  Formulae
\eqref{eq:R}--\eqref{eq:zdef} simplify to
\begin{equation}\label{eq:zexact}
 z_0=\frac{8811237}{43873750}
   +i\frac{4319088\sqrt6}{21936875}.
\end{equation}
If $\phi_0=\arg z_0$, exact arithmetic gives
\begin{equation}\label{eq:inside}
 |z_0|^2=\frac{598702419}{2193687500},\qquad
 R(\phi_0)=\frac{140986860969}{260066040500},
\end{equation}
and hence
\begin{equation}\label{eq:mz}
 m_D(z_0)^2
 =\frac{|z_0|^2}{R(\phi_0)^2}
 =\frac{456889764289102564}{492000073411354561}<1.
\end{equation}
Thus $z_0\in D$.  By construction and \eqref{eq:B},
\[
 B_D(\zeta_0,z_0)
 =s q_0(\zeta_0-z_0)+1-s^2=0.
\]

It remains to check that this is not a harmless isolated coincidence.  Write
$\zeta=te^{i\theta}$ near $\zeta_0$ and regard $B_D(\zeta,z_0)$ as a map
from the real variables $(t,\theta)$ to $\C\simeq\R^2$.  Direct
differentiation of \eqref{eq:B}, \eqref{eq:gauge}, and \eqref{eq:q} at
$(t,\theta)=(sR(\theta_0),\theta_0)$ gives
\begin{align*}
 \partial_t B_D
 &=-\frac{1000}{201}
   +i\frac{26000\sqrt6}{40401},\\
 \partial_\theta B_D
 &=\frac{1825018\sqrt6}{7054899}
   -i\frac{8776}{35099}.
\end{align*}
The corresponding real Jacobian is
\begin{equation}\label{eq:jac}
 \det\frac{\partial(\Rea B_D,\Ima B_D)}{\partial(t,\theta)}
 =\frac{23285456000}{95008324833}>0.
\end{equation}
Since $t=sR(\theta_0)>0$, the change from polar to Cartesian coordinates is
nonsingular, so $D_\zeta B_D(\zeta_0,z_0)$ is invertible as a real map.

The real implicit-function theorem now yields neighborhoods
$z_0\in V\Subset D$ and $\zeta_0\in U\Subset D$ and a $C^1$ map
$z\mapsto\zeta(z)\in U$ such that
\[
 B_D(\zeta(z),z)=0\qquad(z\in V).
\]
After shrinking the neighborhoods, the first derivatives of $B_D$ are
uniformly bounded there.  Therefore, for some $C>0$,
\[
 |B_D(\zeta,z)|\le C|\zeta-\zeta(z)|
\]
whenever $z\in V$ and $\zeta$ lies in a fixed small disk about $\zeta(z)$
contained in $U$.  Taking $f\equiv1\in L^p(D)$ and recalling that $n=1$,
\begin{align*}
 |\mathcal B_D|1(z)
 &\ge C^{-2}\int_{|\zeta-\zeta(z)|<r}
       |\zeta-\zeta(z)|^{-2}\,dA(\zeta)=\infty,
 \qquad z\in V.
\end{align*}
Thus the output is infinite on an open set, proving the theorem.
\end{proof}

\section{Why the real-convex argument breaks}

The source first uses real convexity to show that the real tangent line at a
boundary point does not meet $\overline D$ elsewhere.  That makes the real
part of $B_D$ vanish only on the boundary diagonal.  It then upgrades this to
a two-sided nonisotropic estimate and applies a weighted Schur test.  For
\eqref{eq:R}, the outward normal from one side enters the opposite lobe; the
point \eqref{eq:zdef} converts precisely that geometry into the zero
$B_D(\zeta_0,z_0)=0$.  No Schur estimate can survive an interior transverse
zero of the positive kernel.

\section{Verification, scope, and novelty}

The accompanying exact-arithmetic script independently checks
\eqref{eq:zexact}--\eqref{eq:jac} and the identity $B_D(\zeta_0,z_0)=0$.
No floating-point inequality is used in the proof.  The source PDF is copied
locally and the question crop is taken from page~27.

The scope qualification is essential.  This is a full counterexample to the
quantified question printed in Remark~7.3 because that question does not
exclude $n=1$.  Several-complex-variables authors sometimes reserve the
phrase ``strongly $\C$-convex'' for $n\ge2$, precisely because the complex
tangent condition is empty in dimension one.  Under that intended
reinterpretation, the packet proves only that an $n\ge2$ hypothesis is
necessary; it does not answer the higher-dimensional problem.

A bounded novelty search on 11 August 2026 covered the run's registry,
solution, attempt, and proof-gap indexes; exact-title and exact-question web
searches; arXiv id 2506.02913; and the phrases ``strongly C-convex'',
``$B_D$'', ``Leray kernel'', and ``$L^p$ bounded''.  It found the source,
general Cauchy--Leray literature, and summaries of the source, but no later
paper answering Remark~7.3 and no occurrence of this radial counterexample.
This supports novelty only within those stated bounds.

\paragraph{Human-review recommendation.}
Verify first that the intended convention in Remark~7.3 permits $n=1$.
Conditional on that literal reading, the proof reduces to exact arithmetic
and a standard implicit-function argument and should be checked as a
candidate full negative answer.  If the authors intended $n\ge2$, classify
it instead as a sharp dimensional caveat and retain the higher-dimensional
question as open.

\begin{thebibliography}{9}
\bibitem{Chatterjee}
A. Chatterjee,
\emph{Dual realizations of Bergman spaces on strongly convex domains},
arXiv:2506.02913v2, 2026, Remark~7.3 and Lemma~3.1.
\end{thebibliography}

\end{document}
